RNG builds a datacenter network as a flat random graph of top-of-rack routers, and its routing protocol, Spraypoint, sends every packet through a random neighbor first. That spreads load well, but it sends even a tiny message between two directly connected racks over two to five hops. We test a \emph{bypass} that routes small flows on the shortest path instead. Using an ns-3 simulation of a 256-router fabric, flow-level computations on the paper's 1{,}000-router fabric, and analytical models extended from the RNG paper, we measure how much the bypass speeds up small flows, what it costs the rest of the traffic under uniform and worst-case traffic, and when that cost outweighs the gain.
